\section[The small J-function of Gr {n,N}]{The small J-function of $\mathbf{Gr}_{\boldsymbol{n,N}}$}
\label{Sec2}

{\sloppy
Let $X:={\rm Gr}_{n,N}$ be the grassmannian of $n$-dimensional subspaces $V\subset \CC^N$. Its K-ring $K^0(X)$ is~generated by the exterior powers $\bigwedge^kV$ of the tautological bundle, $k=1,\dots,n$. Using the \mbox{splitting} principle, we will often write them as elementary symmetric functions $\sum_{1\leq i_1<\cdots<i_k\leq n}P_{i_1}\cdots P_{i_k}$ of K-theoretic Chern roots of $V=P_1+\cdots+P_n$.

}

\begin{Proposition} \label{prop1}
The K-theoretic Poincar{\'e} pairing on $K^0(X)$ is given by residue formula
\[
\chi(X; \Phi(P))= (-1)^n\Res_{P=1} \frac{\Phi(P) \ \prod_{i\neq j}(1-P_i/P_j)}{(1-P_1)^N\cdots(1-P_n)^N} \frac{{\rm d}P_1\wedge\cdots \wedge {\rm d}P_n}{P_1\cdots P_n},
\]
where $\Phi$ is any symmetric Laurent polynomial of $P_1,\dots,P_n$.
\end{Proposition}

The formula is obtained as the non-equivariant limit $\Lambda\to 1$ from its $T^N$-equivariant counterpart, where $T^N$ is the torus of diagonal matrices $\operatorname{diag}(\Lambda_1,\dots,\Lambda_N)$ acting on $\C^N$.

\begin{Proposition}\label{prop2}
The $T^N$-equivariant K-theoretic Poincar{\'e} pairing on $K^0_T(X)$ is given by
\[
\chi_T(X;\Phi(P,\Lambda))= \frac{(-1)^n}{n!} \Res_{P\neq 0,\infty} \frac{\Phi(P,\Lambda)\ \prod_{i\neq j}(1-P_i/P_j)}{\prod_{i=1}^n\prod_{j=1}^N(1-P_i/\Lambda_j)} \frac{{\rm d}P_1\wedge\cdots \wedge {\rm d}P_n}{P_1\cdots P_n}.
\]
\end{Proposition}
Here $\Phi$ is a Laurent polynomial in $P$ and $\Lambda$, symmetric in $P$, $\chi_T$ is the $T$-equivariant holomorphic Euler characteristic, taking values in the representation ring $\ZZ[\Lambda^{\pm}]$ of the torus, and the residue sum is taken over all poles $P_1=\Lambda_{i_1}$, \dots, $P_n=\Lambda_{i_n}$ (with this ordering of the equations, and $i_\a\neq i_\b$). The formula is proved by the direct application of Lefschetz' holomorphic fixed point formula.

The following $q$-hypergeometric series has emerged from a study of spaces of rational curves in the grassmannian based on their quasimap compactifications:
\[
J=\sum_{0\leq d_1,\dots,d_n}\frac{Q^{d_1+\cdots+ d_n}}{\prod_{i=1}^n\prod_{m=1}^{d_i}(1-q^mP_i)^N} \prod_{i,j=1}^n\frac{\prod_{m=-\infty}^{d_i-d_j}(1-q^mP_i/P_j)} {\prod_{m=-\infty}^0(1-q^mP_i/P_j)}.
\]

\begin{Remark}\label{rmk1}
The product on the right can be rearranged as
\[
\prod_{d_i>d_j} q^{\binom{d_i-d_j}{2}} \bigg(\!{-}\frac{P_i}{P_j}\bigg)^{d_i-d_j}\frac{P_j-q^{d_i-d_j}P_i}{P_j-P_i},
\]
and therefore may contain nilpotent factors $P_j-P_i$ in the denominator. It is not hard to see, however, that transposing $P_i$ and $P_j$ does not change the sum of terms with a fixed $d_1+\cdots +d_n$, implying that after clearing the denominators, the numerator becomes divisible by $P_j-P_i$ (namely, it changes sign under the transposition, and hence vanishes when $P_i=P_j$).
\end{Remark}

\begin{Remark}\label{rmk2}
Another consequence of the above rearrangement is that, with the exception of the term $Q^0$, the series consists of reduced rational functions of $q$. Namely, the factor at $Q^{d_1+\cdots+d_n}$ has no pole at $q=0$, and the $q$-degree of the denominator exceeds that of the numerator by
\[
N\sum_{i=1}^n\binom{d_i+1}{2}-\sum_{d_i>d_j}\binom{d_i-d_j+1}{2}\geq N-n+1\geq 2.
\]
\end{Remark}

\begin{Theorem}[{cf.~\cite{Taipale, Wen}}]\label{thm1}
The series $(1-q) J$ is the ``small J-function'' of the grassmannian ${\rm Gr}_{n,N}(\CC)$.
\end{Theorem}

Recall that the genus-0 quantum K-theory of a target space $X$, the ``big J-function'' is defi\-ned~as
\[
\t \mapsto \J(\t):=1-q+\t(q)+\sum_{d,m,\a} Q^d \phi_\a \bigg\lan \frac{\phi^\a}{1-qL_0},\t(L_1),\dots,\t(L_m)\bigg\ran_{0,m+1,d}^{S_m}.
\]
Here $Q$ is the Novikov's variable, $\{ \phi_\a \}$ and $\{ \phi^\a\}$ are Poincar{\'e}-dual bases in $K^0(X)$, $\t=\sum_k t_k q^k$ is a Laurent polynomial in $q$ with vector coefficients $t_k \in K^0(X)\otimes \QQ[[Q]]$, and the correlator represents the K-theoretic GW-invariant computes a suitable holomorphic Euler characteristic on the moduli spaces of stable maps $X_{g,m+1,d}:=\M_{g,m+1}(X,d)$ with the input (or ``insertion'') at the marked point (with the index $i=0,\dots,m$ in the above formula) of the form $\sum_k (\ev_i^*t_k) L_i^k$, where $L_i$ stands for the universal cotangent line bundle at the $i$th marked point. From among several flavors of such K-theoretic GW-invariants (ordinary as in~\cite{Givental-Tonita}, or permutation-equivariant as in~\cite{Givental:perm9}), we will currently use the {\em permutation-invariant} ones (as the superscript $S_m$ indicates), i.e., computing the super-dimension of the part of the sheaf cohomology on the moduli space $X_{0,m+1,d}$ which is {\em invariant} under permutations of the $m$ marked points with the indices $i=1,\dots,m$ carrying the symmetric inputs~$\t(L_i)$.

The ``small J-function'' is obtained from $\J$ by setting the input $\t=0$. In particular, this eliminates the role of the permutation group, and so $\J(0)$ represents the ``ordinary'' K-theoretic GW-invariants. Thus, according to Theorem~\ref{thm1},
\[
\J(0):=(1-q)+\sum_{d,\a}Q^d\phi_\a\bigg\lan\frac{\phi^\a}{1-qL_0}\bigg\ran_{0,1,d}=(1-q) J.
\]
Theorem~\ref{thm1} is obtained as the non-equivariant limit $\Lambda\to 1$ from the following result about the $T^N$-equivariant version $\J^T$ ``big J-function'' of the grassmannian.

\begin{Theorem}[{cf.~\cite{Taipale, Wen}}]\label{thm2}
$\J^T(0)=(1-q)J^T$, where
\[
J^T=\sum_{0\leq d_1,\dots,d_n}\frac{Q^{d_1+\cdots+d_n}}{\prod_{i=1}^n\prod_{j=1}^N\prod_{m=1}^{d_i}(1-q^mP_i/\Lambda_j)} \prod_{i,j=1}^n\frac{\prod_{m=-\infty}^{d_i-d_j}(1-q^mP_i/P_j)} {\prod_{m=-\infty}^0(1-q^mP_i/P_j)}.
\]
\end{Theorem}

Note that by the very definition (the same as for $\J$ with the correlators taking values in the representation ring $\ZZ[\Lambda^{\pm}]$), the function $\J^T$ is a $Q$-series with coefficients which are rational functions of $q$ with vector values in $K:=K^0_T(X)\otimes \QQ[[Q]]$. Abusing the language we call such series rational functions of $q$, denote the space they form by $\K:=K(q^{\pm})$, and
call it the {\em loop space}. The part $(1-q)+\t(q)$ (``dilaton shift''$+$''input'') belongs to the subspace $\K_{+}$ consisting of Laurent polynomials (they can have poles only at $q=0,\infty$), while the sum of the correlators belongs (as it is not hard to see) to the complementary subspace $\K_{-}=\{ \f \in \K\ | \ \f(0)\neq \infty, \f(\infty)=0.\}$. It follows from Remark~\ref{rmk2} above (which applies to $J^T$ as well) that $(1-q)J^T\equiv 1-q\ \mod \K_{-}$. Thus, the non-obvious statement of Theorem~\ref{thm2} is that $(1-q)J^T$ represents a value of $\J^T$ at all.

The technique of fixed point localization we intend to use goes back to paper~\cite{Jeff_Brown} by
J.~Brown, and was adapted to the K-theoretic situation in~\cite{Givental:perm2}. The technique, applicable whenever the target carries a torus action with isolated fixed points and isolated $1$-dimensional orbits, completely characterizes all values of the big J-function $\J^T$ as the set of those rational functions $\f\in \K$ which pass two tests: criterions (i) and (ii). They are formulated in terms of specializations $\f_{\a}$ of $\f$ to the fixed point of the torus action in $X$. In the case of the grassmannian, take for example the fixed point $V_{1,\dots,n}=\mathop{\rm Span}(e_1,\dots,e_n)$ where (we may assume by choosing the ordering) $(P_1,\dots,P_n)=(\Lambda_1,\dots,\Lambda_n)$:
\[
J^T_{(1,\dots,n)}= \sum_{0\leq d_1,\dots,d_n} \frac{Q^{d_1+\cdots+d_n}}{\prod_{i=1}^n\prod_{j=1}^N\prod_{m=1}^{d_i}(1-q^m\Lambda_i/\Lambda_j)} \prod_{i,j=1}^n\frac{\prod_{m=-\infty}^{d_i-d_j}(1-q^m\Lambda_i/\Lambda_j)} {\prod_{m=-\infty}^0(1-q^m\Lambda_i/\Lambda_j)}.
\]
Note that the 1st factor contains the product (coming from $j=i$):
\[
\frac{1}{\prod_{i=1}^n \prod_{m=1}^{d_i}(1-q^m)}
\]
with poles at roots of unity, while all other poles are elsewhere (at $q=(\Lambda_i/\Lambda_j)^{-1/m}$). Each term of the series considered as rational functions of $q$ can be split (e.g., using partial fraction decomposition) into the sum of a reduced rational function with poles at the roots of unity and a rational function with poles elsewhere. The result will be interpreted (or rather termed) as~a~mero\-mor\-phic function in a neighborhood of the roots of unity.

Criterion (i) stipulates that {\em $\f_\a$, when interpreted as a meromorphic function in a neighborhood of the roots of unity, must represent a value (over a suitable ground ring) of the big J-function of the point target space.} We will return to this criterion in the next section and explain how it~can be verified.

Criterion (ii) controls residues of $\f_\a(q){\rm d}q/q$ at the poles originating from $T$-equivariant covers of $1$-dimensional orbits. Namely, the tangent space to the grassmannian at the fixed point $V_{(1,\dots,n)}$ carries the torus action with the distinct eigenvalues $\Lambda_j/\Lambda_i$, $i=1,\dots,n$, $j=n+1,\dots,N$. Consequently, for each choice of $i$ and $j$ there is a $1$-dimensional orbit, which
compactifies into~$\CC P^1$ connecting this fixed point with another one. For instance, taking $i=1$ and $j=n+1$, we find such an orbit connecting $V_{(1,\dots,n)}$ with $V_{(2,\dots,n+1)}$. Let $\phi\colon \CC P^1\to \CC P^1$ be the map $z\mapsto z^{m_0}$ ramified at $z=0,\infty$ (representing the two fixed points which we call $\a$ and $\b$). Criterion (ii) has the form of the recursion relation:
\[
\Res_{q=(\Lambda_j/\Lambda_i)^{1/m_0}}\f_\a(q) \frac{{\rm d}q}{q} = -\frac{Q^{m_0}}{m_0}
\frac{\Eu(T_\a X)}{\Eu(T_{\phi}X_{0,2,m_0})} \, \f_{\b}\Big|_{q=(\Lambda_j/\Lambda_i)^{1/m_0}},
\]
where $\Eu$ are equivariant K-theoretic Euler classes: of the tangent space to $X$ at $\a$, and to the moduli space of degree-$m_0$ stable maps with 2 marked points at the point represented by the $m_0$-fold cover $\phi$ respectively.

We compute $\Res_{q=(\Lambda_{n+1}/\Lambda_1)^{1/m_0}} (1-q)J^T_{(1,\dots,n)}(q)dq/q$, replacing $d_1$ with $d_1+m_0$, assuming $d_i=d_1$ when $i=n+1$, and using $x:= (\Lambda_{n+1}/\Lambda_1)^{1/m_0}$:
\begin{gather*}
\Res_{q=x} (1-q)J^T_{(1,\dots,n)}(q)\frac{{\rm d}q}{q} = -(1-x)\frac{Q^{m_0}}{m_0}
\\ \qquad
{}\times\frac{1}{\prod_{j=1}^N\prod_{m=1}^{m_0}\,_{(j,m)\neq (n+1,m_0)}\ (1-x^m\Lambda_1/\Lambda_j)} \prod_{j=2}^n\prod_{m=1}^{m_0}\frac{1-x^m\Lambda_1/\Lambda_j}{1-x^m\Lambda_j/\Lambda_{n+1}}
\\ \qquad
{}\times\sum_{0\leq d_1,\dots,d_n}\frac{Q^{d_1+\cdots+d_n}}{\prod_{i=2}^{n+1}\prod_{j=1}^N\prod_{m=1}^{d_i}
(1-x^m\Lambda_i/\Lambda_j)} \prod_{i,j=2}^{n+1}\frac{\prod_{m=-\infty}^{d_i-d_j}
(1-x^m\Lambda_i/\Lambda_j)}{\prod_{m=-\infty}^0(1-x^m\Lambda_i/\Lambda_j)}.
\end{gather*}
The sum together with the factor $1-x$ yields $(1-q)J^T_{(2,\dots,n+1)}(q)|_{q=x}$, so it remains to interpret the recursion coefficient in terms of the Euler classes.

Applying Lefschetz' fixed point formula on ${\rm Gr}_{n,N}(\CC)$, we've already used that $\Eu(T_{(1,\dots,n)}X)\!=\prod_{i=1}^n\prod_{j=n+1}^N(1-\Lambda_i/\Lambda_j)$. In order to compute $\Eu(T_{\phi}X_{0,2,m_0})$, we note that the $1$-dimensional orbit connecting the fixed points $\mathop{\rm Span} (e_1,\dots,e_n)$ and $\mathop{\rm Span} (e_2,\dots,e_{n+1})$ consists of subspaces $V_t:=\mathop{\rm Span}(t e_1+(1-t) e_{n+1}, e_2,\dots,e_{n-1})$. Consequently, restricted to $\CC P^1=\{ V_t \}$, the tangent bundle to the grassmannian, which has the form $\mathop{\rm Hom} (V_t,\CC^N/V_t)$, can be described in terms of~the~Hopf bundle $L$ over $\CC P^1$ and its complementary $L':=\mathop{\rm Span}(e_1,e_{n+1})/L=\Lambda_1\Lambda_{n+1}L^{-1}$ as
\[
\mathop{\rm Hom} \big(L \oplus \mathop{\rm Span}(e_2,\dots,e_n), L' \oplus \mathop{\rm Span}(e_{n+2},\dots,e_N)\big).
\]
On the $m_0$-fold cover $\phi\colon \CC P^1\to \CC P^1$, the contributions to the Euler class of the $T$-modules $H^0\big(\CC P^1; \phi^*L^{-1}\otimes \mathop{\rm Span} (e_j)\big)$ and $H^0\big(\CC P^1; \phi^*L'\otimes \mathop{\rm Span}(e_i)^{-1}\big)$ are respectively
\[
\prod_{m=0}^{m_0}(1-x^m\Lambda_1/\Lambda_j)\qquad \text{and}\qquad \prod_{m=0}^{m_0}(1-x^m\Lambda_i/\Lambda_{n+1}).
\]
The contribution of $H^0\big(\CC P^1; \phi^*L^{-1}\otimes L'\big)$ is (as in the case of $X=\CC P^1$) $\prod_{m=-m_0, m\neq 0}^{m_0}(1-x^m)$. Combing all the contributing factors, we find
\begin{gather*}
\frac{\Eu(T_{(1,\dots,n)}X)}{\Eu(T_{\phi}X_{0,2,m_0})}=
\frac{\prod_{i=1}^n\prod_{j=n+1}^N(1-\Lambda_i/\Lambda_j)}
{\prod_{i=2}^n\prod_{j=n+2}^N(1-\Lambda_i/\Lambda_j)}
\frac{1}{\prod_{m=-m_0, m\neq 0}^{m_0}(1-x^m)}
\\ \hphantom{\frac{\Eu(T_{(1,\dots,n)}X)}{\Eu(T_{\phi}X_{0,2,m_0})}=}
{}\times \frac{1}{\prod_{j=n+2}^N\prod_{m=0}^{m_0}(1-x^m\Lambda_1/\Lambda_j)}
\frac{1}{\prod_{i=2}^n\prod_{m=0}^{m_0}(1-x^m\Lambda_i/\Lambda_{n+1})}.
\end{gather*}
Checking that this expression matches exactly the recursion coefficient for $J^T$ (the middle line in~the formula for the residue) is the matter of a straightforward (though somewhat cumbersome) rearrangement of the factors.

\begin{Remark}\label{rmk3}
Note that the structure of the recursion relations (ii) and the values of the recursion coefficients completely characterize the big J-function of a particular theory, since criterion (i) does not involve any additional choices.
\end{Remark}

\section{Quantum K-theory of the point}
\label{Sec3}

Genus-0 permutation-equivariant K-theoretic GW-invariants of the point are represented by the ``big J-function'' of the form
\[
\J_{pt}(\t):=(1-q)+\t(q)+\sum_{m=2}^{\infty}\chi
\bigg(\M_{0,m+1}/S_m; \frac{1}{1-L_0q}\otimes_{i=1}^m \t(L_i)\bigg).
\]
Here $L_i$ are the universal cotangent line bundles over the Deligne--Mumford spaces $\M_{0,m+1}$. The~holomorphic Euler characteristic $\chi$ on the orbifold $\M_{0,m+1}/S_m$ computes the super-dimen\-sion of the $S_m$-invariant part of the sheaf cohomology on $\M_{0,m+1}$. The insertions $\t(L_i)$ (and hence the input $\t(q)$) can be in fact any rational functions of $L_i$ (respectively of $q$) as long as they don't have poles at roots of unity. On the contrary, each $\chi$-term is a reduced rational function of $q$ with poles only at the roots of unity (of order $\leq m$). Both this and the previous claim easily follow from the general structure of the Lefschetz fixed point formula (applied on $\M_{0,m+1}/S_m$).

As it is explained in~\cite{Givental:perm1, Givental:perm9}, the action of permutation groups on the sheaf cohomology is captured by the above $S_m$-invariants taking values in an arbitrary $\lambda$-algebra, i.e., a ring $R$ equipped with the Adams operations $\Psi^k\colon R\to R$. E.g., for $\tau\in R$
\[
\chi\big(\M_{0,m}/S_m;\otimes_{i=1}^m \big(\tau L_i^d\big)\big) := \frac{1}{m!} \sum_{h\in S_m}\prod_{k=1}^{\infty} (\Psi^k\tau)^{l_k(h)} \str_h H^*\big(\M_{0,m};\otimes_{i=1}^mL_i^d\big),
\]
where $l_k(h)$ denotes the number of cycles of length $r$ in the cycle decomposition of permutation $h$. For mode detail, we refer to~\cite{Givental:perm1} or~\cite{Givental:perm9}, where this example is
extrapolated to general $R$-valued insertions $\t(L_i)$. Note that only $\Psi^r$ with $r>0$ are used in this definition.

We assume that $R$ is complete in the adic topology defined by a certain ideal $R_{+}\subset R$, which is respected by the Adams operations $\Psi^k$ with $k>0$ in the sense that $\Psi^k(R_{+})\subset R_{+}^k$, and that the input $\t$ is ``small'' in the sense that it takes values in $R_{+}$. In fact the latter property guarantees that the $m$-th $\chi$-term of the series $\J_{pt}$ takes values in $R_{+}^m$, with assures $R_{+}$-adic convergence of the series.

The genus-$0$ permutation-equivariant GW-invariants of the point target space are completely described in~\cite{Givental:perm3}. Namely, given a ground $\lambda$-algebra $R$, the range of the big J-function
$\t\mapsto \J_{pt}(\t)$, which is a semi-infinite cone (that we will denote $\L_{pt}$) in the (completed) space of $R$-valued rational functions of $q$ (which we will denote $R(q^{\pm})$) is explicitly parameterized as
\[
(1-q) {\rm e}^{ \sum_{k>0}\Psi^k(\tau)/k(1-q^k)}\big(1+R_{+}[q^{\pm}]\big).
\]
Here $\tau\in R_{+}$, and the notation $R_{+}[q^{\pm}]$ is reserved for the completion in the $R_{+}$-adic topology of~the~space of rational functions of~$q$ which have no poles at roots of unity and take values in~$R_{+}$.

In our arguments, we will take advantage of the possibility to replace one ground $\lambda$-algebra with another related to it by a homomorphism respecting the Adams operations. In simple terms: If some
$\f\in R(q^{\pm})$ is known to lie in $\L_{pt}$ (over a given ground ring~$R$), i.e., the part $\f_{-}$ with poles at the roots of unity represents K-theoretic GW-invariants with the input defined by~the part $\f_{+}$ with poles away from roots of unity, the same will remain true when the values of some parameters (coordinates on $\operatorname{Spec} R$) are specialized in a way commuting with $\Psi^k$ for all~$k>0$.

Another important property of the cone $\L_{pt}$ that we will rely on is its invariance under a~cer\-tain group of (pseudo) finite-difference operators. Namely, let $R=\QQ[[Q]]$, where $\Psi^kQ=Q^{|k|}$, $k=\pm 1,\pm 2,\dots$, and let $D(q^{Q\p_Q},Q)$ be a finite difference operator (which we should assume ``small'' in $R_{+}$-adic sense to assure convergence). It is almost obvious that the linear vector field $\f \mapsto D\f$ in $R(q^{\pm})$ is tangent to $\L_{pt}$, and therefore ${\rm e}^D$ preserves $\L_{pt}$. Moreover, according to a~result from~\cite{Givental:perm4}, $\L_{pt}$ is preserved by the operator
\[
\f\mapsto {\rm e}^{ \sum_{k>0} \Psi^k(D(q^{kQ\p_Q},Q))/k(1-q^k)} \f.
\]

Our goal in this section is to verify that the series $J^T_{(1,\dots,n)}$ from the previous section satis\-fies criterion~(i), i.e., that it represents a value of $\J_{pt}$ when interpreted as a meromorphic function of $q$ in a neighborhood of roots of unity. For this, we begin with the ground ring $R=\QQ\big[\Lambda_1^{\pm},\dots,\Lambda_N^{\pm}\big][[Q_1,\dots,Q_n]]$, with the Adams operations acting by $\Psi^k\big(\Lambda_j^{\pm 1}\big)=\Lambda_j^{\pm k}$, $\Psi^k(Q_i)=Q_i^{|k|}$, and set $R_{+}=(Q_1,\dots,Q_n)$. In particular, taking $\tau=Q_1+\cdots+Q_n$, we~find that $\L_{pt}\ni (1-q) J_{pt}$, where $J_{pt}$ is the following product of $q$-exponential functions:
\begin{gather*}
J_{pt}:= {\rm e}^{\sum_{i=1}^n\sum_{k>0}Q_i^k/k(1-q^k)}
= \prod_{i=1}^n \sum_{d_i=0}^{\infty}\frac{Q_i^{d_i}}{\prod_{m=1}^{d_i}(1-q^m)}
 = \sum_{0\leq d_1,\dots,d_n}\frac{Q_1^{d_1}\cdots Q_n^{d_n}}{\prod_{i=1}^n\prod_{m=1}^{d_i}(1-q^m)}.
 \end{gather*}

Following~\cite{Givental:perm4}, we are going to use $Q$-independent finite-difference operators of the form $D_{l,\Lambda}:=-q \Lambda (1-q^{l\cdot Q\p_Q})$, where $l\cdot Q\p_Q:=\sum_i l_iQ_i\p_{Q_i}$, and $\Lambda$ is a formal variable added to the ground ring, $\Psi^k\Lambda:=\Lambda^{|k|}$. The corresponding $\L_{pt}$-preserving operator is
\[
\Gamma_{l,\Lambda}:= {\rm e}^{ -\sum_{k>0}\Lambda^k
(1-q^{kl\cdot Q\p_Q})q^k/k(1-q^k)}.
\]
This expression is in fact the asymptotical expansion (near the unit circle on the $q$-plane) of the ratio:
\[
\prod_{m=-\infty}^0\big(1-\Lambda q^{l\cdot Q\p_Q}q^m\big)/ \prod_{m=-\infty}^0(1-\Lambda q^m).
\]
The ratio and its asymptotical expansion act on monomials $Q^d=Q_1^{d_1}\cdots Q_n^{d_n}$ the same way:
\[
\Gamma_{l,\Lambda}Q^d= Q^d\ \frac{\prod_{m=-\infty}^0(1-\Lambda q^{m+(l\cdot d)})}{\prod_{m=-\infty}^0(1-\Lambda q^m)} = Q^d\ \frac{\prod_{m=-\infty}^{l\cdot d}(1-\Lambda q^m)}{\prod_{m=-\infty}^0(1-\Lambda q^m)}.
\]
Note that the right hand side is a rational function of $q$ with poles away from roots of unity. Considering $\Gamma_{l,\Lambda}(1-q)J_{pt}$ as a point of $\L_{pt}$ over the ground ring $R[[\Lambda]]$, we conclude therefore, that being a $Q$-series with coefficients which are rational function of $q$ with coefficients polynomial in $\Lambda$, it is a point of $\L_{pt}$ over $R[\Lambda]$. Thus, it will turn into a point of $\L_{pt}$ over $R$ when $\Lambda$ is replaced by any non-trivial monomial from $\QQ[\Lambda_1^{\pm},\dots,\Lambda_N^{\pm}]$, e.g., by $\Lambda_i/\Lambda_j$ with $i\neq j$.

In order to complete our fixed point localization proof of Theorem~\ref{thm2}, we apply to $J_{pt}$ the following operators (where $l=\1_i$ contains $1$ in the $i$-th position and $0$ everywhere else):
\begin{gather*}
\bigg(\prod_{i=1}^n\prod_{j=1, j\neq i}^N \Gamma^{-1}_{\1_i,\Lambda_i/\Lambda_j}\bigg) \bigg(\prod_{i,j=1}^n\Gamma_{\1_i-\1_j,\Lambda_i/\Lambda_j}\bigg)\ J_{pt}
\\ \qquad
{}=\sum_{0<d_1,\dots,d_n}\frac{Q_1^{d_1}\cdots Q_n^{d_n}}{\prod_{i=1}^n\prod_{j=1}^N\prod_{m=1}^{d_i}(1-q^m\Lambda_i/\Lambda_j)} \prod_{i,j=1}^n\frac{\prod_{m=-\infty}^{d_i-d_j}(1-q^m\Lambda_i/\Lambda_j)}{\prod_{m=-\infty}^0(1-q^m\Lambda_i/\Lambda_j)}.
\end{gather*}
The terms of the last sum are interpreted as meromorphic functions in the neighborhood of~roots of unity, i.e., with poles (which come from the factors with $i=j$ in the left product) at the roots of unity only.

When multiplied by $1-q$, the latter series lies in $\L_{pt}$ over the ground $\lambda$-algebra $R=\QQ[\Lambda^{\pm}][[Q_1,\dots,Q_n]]$. The substitution $Q_1=\cdots=Q_n=Q$ (which induces a homomorphism of~$\lambda$-algebras $R\to R_0:=\QQ[\Lambda^{\pm}][[Q]]$) yields therefore a series which lies in the range $\J_{pt}$ over~$R_0$. It actually coincides with the localization $(1-q) J^T_{(1,\dots,n)}$ of $J^T$ at the indicated fixed point. Therefore $(1-q) J^T_{(1,\dots,n)}$, when interpreted as a series of meromorphic functions near roots of~unity, satisfies criterion (i). Due to the Weyl group symmetry between all fixed points, and between all $1$-dimensional orbits connecting them, this finishes the proof of Theorem~\ref{thm2}.


